% !PRINTED 310
\weberpage{310}
%% [page furniture, not body text -- commented out 2026-08-10] Page furniture (running
%% head, folio, signature) is excluded from this edition by standing policy.
%% \noindent Zehnter Abschnitt. \hfill \S.\,81.

%<<CONT
and since the system of the $z_i'$ runs through the same sequence of values as the system of the $z_i$, this sum is equal to $\Psi_{\xi_1',\xi_2',\xi_3'}$, and we obtain:

\gesperrt{2. The application of the permutation $\sigma$ to the $X$ in the functions $\Psi$ is equivalent to the application of $\varrho^{-1}$ to the indices $\xi_1, \xi_2, \xi_3$.}

\begin{center}
\S.\,81.\\
\gesperrt{The triple systems of the resolvents.}
\end{center}

We now denote the resolvent $\Psi_{\xi_1,\xi_2,\xi_3}$ more briefly by $\Psi^\xi$ and remark that, by formula (2), \S.\,80, besides the squares of these functions certain products also admit the permutations of the group $Q$. To these belongs first of all the product of all the quantities $\Psi$, which we shall denote by $A$:
\weq{1}{A = \overset{\xi}{\prod} \Psi_\xi,}
but then also the products of three, $\Psi_\xi \Psi_\eta \Psi_\zeta$, and consequently also their reciprocal values
\weq{2}{v = \frac{1}{\Psi_\xi \Psi_\eta \Psi_\zeta},}
when the indices satisfy the conditions:
\weq{3}{\begin{aligned}
\xi_1 + \eta_1 + \zeta_1 &\equiv 0\\
\xi_2 + \eta_2 + \zeta_2 &\equiv 0 \quad (\mathrm{mod}\ 2).\\
\xi_3 + \eta_3 + \zeta_3 &\equiv 0
\end{aligned}}
Such a system of three functions $\Psi$ we call a \gesperrt{triple}. Likewise we shall call three index systems $(\xi), (\eta), (\zeta)$ that satisfy the conditions (3) a triple.

There are in all seven and no more such triples; for among the three systems $(\xi), (\eta), (\zeta)$ no two can be equal to one another, and by two of them the third is uniquely determined. One can therefore take for $(\xi), (\eta)$ any pair from the seven possible systems $(\xi)$, but then obtains each triple six times. The total number is therefore $7 . 6 : 6 = 7$.

In order to characterise the individual triple systems, we introduce three numbers $\alpha_1, \alpha_2, \alpha_3$, to be taken modulo 2, which do not all three vanish, and remark that the congruence
\weq{4}{\alpha_1 \xi_1 + \alpha_2 \xi_2 + \alpha_3 \xi_3 \equiv 0 \ (\mathrm{mod}\ 2)}
%CONT>>
